\section{The actual second radial coefficient and its regular arithmetic value}
\label{sec:ns-actual-second-radial}

We now evaluate the regular Mellin datum in
\eqref{ndm:arithmetic-balance} for the complete source field, not an
auxiliary test function. The calculation works along the inner axis,
and retains both powers of the concentration coordinate and the exact
force remainder. The source existence theorem is still imported from
\cite{openaiNS2026release}; the coefficient extraction and its arithmetic
image below are consequences proved here. All coefficients are ordinary
Taylor coefficients in the physical variable \(\sigma=r^2/2\).

\subsection{Source axis data, coordinate inverse, and the full local field}

Retain the source's fixed choices
\[
 0<h<1/100,\quad A=\tfrac12+h,\quad D=\tfrac12-h,\quad
 k=1+h,\quad C>1,\quad 0<j_0\leq1/20,\quad
 \Lambda\geq1,\quad \sigma_*>0 .
\]
These are the constants chosen by the source construction, with its
thresholds and dependencies, not independently variable parameters
asserted to produce solutions. Here \(\sigma_*\) is the source's
positive regularizing parameter, not the physical radial coordinate.
At physical axial coordinate \(z\), put \(y=z/\sqrt\nu\) and retain
\begin{equation}\label{nsr:coordinates}
 \tau=1-t=q(1-\eta^2),\qquad y=q^D\eta,\qquad
 X=\frac{\sigma}{\nu q},\qquad
 d=1-\eta^2,\quad L=1-2h\eta^2.
\end{equation}
The inverse solves \(q-y^2q^{2h}=\tau\), uniquely on
\(q>|y|^{1/D}\), and then sets \(\eta=yq^{-D}\).
The left side starts at zero, tends to infinity, and has derivative
\(1-2hy^2q^{2h-1}=L\geq1-2h>0\) there. Its inverse is smooth
for \(\tau>0\), and \(|\eta|<1\). Thus this is the original
coordinate map, including \(\nu\), with a proved inverse.

Source (B.1), (B.3), and (B.12), pp.~144--147, specify
\begin{equation}\label{nsr:profiles}
 \begin{split}
 U(\eta)&=4\eta+j_0,\qquad H_*(\eta)=D\eta+dU(\eta),\\
 \zeta_*(\eta)&=-\frac{L H_*}{H_*^2+\sigma_*^2},\qquad
 \varphi(\eta)=\exp\left(\Lambda\int_0^\eta\zeta_*(w)\,dw\right).
 \end{split}
\end{equation}
The function \(\zeta_*\) is the source auxiliary function, not the
Riemann zeta function. In particular \(\varphi>0\) on the real
interval and \(\varphi'=\Lambda\zeta_*\varphi\).

The actual axis traces, on the inner localized domain, are
\begin{equation}\label{nsr:axis-traces}
 w_0(z,t)=(u_\nu)_3(0,0,z,t)=\sqrt\nu q^{-A}U(\eta),\qquad
 b_1(z,t)=\partial_\sigma B(0,z,t)=\frac2C q^{-k}\varphi(\eta).
\end{equation}
Here is why these are exact traces of the complete field. Source
Lemma~5.1, p.~47, gives \(\varphi_n(0,\eta)=U_n(0,\eta)=0\)
at every positive base order. Smoothness factors each such profile
by \(X\), and its streamfunction primitive by \(X^2\).
The base cutoffs \(\chi(c_nq)\) are applied to streamfunctions
before taking their curls. Their extra radial term in (5.45), p.~61,
is also divisible by \(X^2\) in \(r u_r\); consequently it
contributes only \(O(r^3)\) to \(u_r\). The axial corrections are
\(O(r^2)\) and swirl corrections \(O(r^3)\). Every fixed
\(z,t\) derivative of these vanishing axis traces is zero as well.
At each preterminal point only finitely many cutoff terms are
present locally, since \(q>0\) and the cutoff scales increase to
infinity. The annular wave and mean representatives, including finite
initialization, vanish in an axis neighborhood by Proposition~9.9,
pp.~114--117. Their locally finite sum has the same property.
Finally (10.4)--(10.5), p.~118, have localization cutoffs identically
one in a fixed neighborhood of the concentrating origin at sufficiently
late times. Their derivative terms vanish there. These facts prove
\eqref{nsr:axis-traces} on an open axis neighborhood, not just at one
point, and justify its axial and temporal derivatives. The physical
map \(u_\nu(x,t)=\sqrt\nu u_*(x/\sqrt\nu,t)\), source (10.22),
cancels the velocity factor against the radial factor in
\(u_\theta/r\), giving precisely the second trace above.

This argument does not discard positive-order corrections to \(b_2\).
Their radial derivatives need not vanish. We determine their total
contribution using the full momentum equation.

\subsection{Exact coefficient extraction, with the force torque retained}

Locally at the axis, the remaining smooth axisymmetric field has the
Cartesian expression
\[
 u_1=a(r^2,z,t)x_1-b(r^2,z,t)x_2,\quad
 u_2=b(r^2,z,t)x_1+a(r^2,z,t)x_2,\quad
 u_3=w(r^2,z,t).
\]
At \(r=0\), incompressibility gives \(2a_0+\partial_z w_0=0\),
and \(b_1=2b_0\). The original angular averages therefore give
\[
 B=b_1\sigma+b_2\sigma^2+O(\sigma^3),\quad
 P=p_2\sigma^2+O(\sigma^3),\quad Q=q_1\sigma+O(\sigma^2),
\]
with the exact coefficients
\begin{equation}\label{nsr:fluxes}
 p_2=4a_0b_0=-b_1\partial_z w_0,\qquad q_1=w_0b_1.
\end{equation}
For example \(P=r^2(ar)(br)=4ab\sigma^2\), which verifies
the factor four without changing the radial variable.

Set
\(r_f(z,t)=(\nabla_x\times f_\nu)_3(0,0,z,t)\).
For any smooth Cartesian force, expanding to first order transversely
and averaging \(x_1f_2-x_2f_1\) gives
\[
 F(\sigma,z,t)=r_f(z,t)\sigma+O(\sigma^2).
\]
Indeed \(\langle x_1^2\rangle=\langle x_2^2\rangle=\sigma\)
and \(\langle x_1x_2\rangle=0\); the surviving coefficient is
\(\partial_1f_2-\partial_2f_1\). This argument does not assume an
axisymmetric force away from the axis.

Take the coefficient of \(\sigma\) in the complete equation
\eqref{nam:pde}. Taylor's differentiated remainder and preterminal
smoothness justify the operation. Substituting \eqref{nsr:fluxes}
in \eqref{ndm:axis-balance} proves
\begin{equation}\label{nsr:extraction}
 \boxed{\quad
 b_2=\frac{\partial_t b_1+w_0\partial_z b_1
          -b_1\partial_z w_0-\nu\partial_{zz}b_1-r_f}{4\nu}.
 \quad}
\end{equation}
Pressure is absent because its full angular torque integrates to
zero, as proved in \eqref{nam:cartesian}--\eqref{nam:pde}, not
because the source pressure or force has been removed.

\subsection{The complete coefficient along the inner axis}

For a function \(v=v(\eta)\), define the explicit operators
\[
 T_bv=\frac{-bv+D\eta v'}L,\qquad
 Z_bv=\frac{2b\eta v+dv'}L.
\]
Implicit differentiation of \eqref{nsr:coordinates} gives
\[
 q_t=-L^{-1},\quad \eta_t=D\eta/(qL),\qquad
 q_z=\frac{2\eta q^{1-D}}{\sqrt\nu L},\quad
 \eta_z=\frac{d q^{-D}}{\sqrt\nu L}.
\]
Consequently, with \(z,t\) physical and the other physical variable
held fixed in each derivative,
\begin{equation}\label{nsr:operators}
 \partial_t(q^bv)=q^{b-1}T_bv,\quad
 \partial_z(q^bv)=\nu^{-1/2}q^{b-D}Z_bv,\quad
 \partial_{zz}(q^bv)=\nu^{-1}q^{b-2D}Z_{b-D}Z_bv.
\end{equation}
This is source (4.2), pp.~24--25, at the axis, derived with the
physical viscosity map. In particular the order of the last two
operators and their different subscripts are part of the formula.

Put
\begin{equation}\label{nsr:HR}
 \begin{split}
 H(\eta)&=T_{-k}\varphi+U Z_{-k}\varphi-\varphi Z_{-A}U,\\
 R(\eta)&=Z_{-k-D}Z_{-k}\varphi.
 \end{split}
\end{equation}
These are completely specified functions, not uncomputed source
remainders. If desired, the second is a rational expression times
the original exponential: writing \(V=-2k\eta\varphi+d\varphi'\),
\[
 R=\frac{2(-k-D)\eta V+d V'}{L^2}-\frac{d V L'}{L^3},
 \quad V'=-2k\varphi-2(k+1)\eta\varphi'+d\varphi'',
 \quad \varphi''=(\Lambda\zeta_*'+\Lambda^2\zeta_*^2)\varphi.
\]
All denominators are nonzero on \([-1,1]\), and the given expressions
specify every derivative appearing in them.
Using \(A+D=1\), substitution in \eqref{nsr:extraction} gives
\begin{equation}\label{nsr:full-b2}
 \boxed{\quad
 b_2(z,t)=\frac{q^{-2-h}H(\eta)-q^{-2+h}R(\eta)}{2\nu C}
                         -\frac{r_f(z,t)}{4\nu}.
 \quad}
\end{equation}
This identity holds throughout the inner axis domain just specified.
In deriving its first exponent, the two advective terms each have
power \(q^{-k-A-D}=q^{-k-1}\). The axial diffusion term has
power \(q^{-k-2D}=q^{-2+h}\). Neither power has been replaced
by the other, and \(r_f\) is the actual final force curl.

The first profile has a useful exact sign. Expanding the numerator
in \eqref{nsr:HR}, using \(U'=4\), \(2(A-k)=-1\), and then
\(\varphi'=\Lambda\zeta_*\varphi\), proves
\begin{equation}\label{nsr:H-sign}
 H(\eta)=\varphi(\eta)
 \left\{\frac{h-3-j_0\eta}{L}
                  -\frac{\Lambda H_*(\eta)^2}{H_*(\eta)^2+\sigma_*^2}\right\}<0
                    \qquad(-1\leq\eta\leq1).
\end{equation}
For the expansion before substitution, the coefficient of \(\varphi\)
is \(k-\eta U-4d=h-3-j_0\eta\), and the coefficient of
\(\varphi'\) is \(D\eta+dU=H_*\), both divided by \(L\).
The strict sign follows from \(\varphi>0\), \(L>0\),
\(h-3+j_0<0\), and the nonnegative square in the second term.
It does not require \(H_*\ne0\).

\subsection{A finite radial closure inside the full cutoff sum}

The source recursion gives more than flatness for this particular
force coefficient. Denote the actual positive-order base cutoffs by
\(\chi(c_nq)\), and retain the profiles \(\varphi_n(X,\eta)\)
from (5.1). The original recurrence (5.3), p.~46, has
\(b=-k\), \(\lambda_n=2nh\), and its full operators include
the \(X\partial_X\) terms in source (4.2). At \(X=0\) it gives
\begin{equation}\label{nsr:finite-radial-closure}
 4\partial_X\varphi_0(0,\eta)=H(\eta),\qquad
 4\partial_X\varphi_1(0,\eta)=-R(\eta),\qquad
 \partial_X\varphi_n(0,\eta)=0\quad(n\geq2).
\end{equation}
Here is the complete coefficient argument. At positive order the
time operator vanishes on the identically zero axis trace
\(\varphi_n(0,\eta)\). In each axial transport product
\(U_i Z_{b+\lambda_j}\varphi_j\), either \(i>0\), in which
case \(U_i(0,\eta)=0\), or \(j>0\), in which case both
the trace and its \(\eta\) derivative vanish. Its
\(X\partial_X\) term also vanishes at zero. For radial transport,
\(V_0=O(X)\), while \(V_i=O(X^2)\) for \(i>0\): the latter
follows by integrating \(\partial_XV_i=-Z_{-A+\lambda_i}U_i\)
and using \(U_i=O(X)\). If \(i=0,j>0\), the factor
\(\varphi_j'+\varphi_j/X\) is bounded, so its product with
\(V_0\) tends to zero. If \(i>0,j=0\), that factor is
\(O(X^{-1})\) but \(V_i=O(X^2)\), again giving zero.
The case with both indices positive also vanishes.

Thus only the negative axial-viscosity term
\(-Z_{b+\lambda_{n-1}-D}Z_{b+\lambda_{n-1}}\varphi_{n-1}\)
can survive. Restriction of each full operator to \(X=0\)
equals its axis operator \(Z\) used above: applying the second
operator does not revive an \(X\partial_X\) term, since it
still contains a factor \(X\) multiplying a smooth function.
The term is \(-R\) for \(n=1\) and zero for every \(n\geq2\).
The left side of (5.3) is \(4\varphi_{n,X}\) at zero,
proving the positive-order assertions.

At leading order the angular equation with leading residual stress
zero, (4.13), has no lower-order axial viscosity. Alternatively its
angular transport expression is the same as (5.3) with indices zero
and that viscosity term omitted. Incompressibility gives
\(\lim_{X\downarrow0}V_0/X=-Z_{-A}U\). Its radial transport
limit is therefore \(-\varphi Z_{-A}U\). Together with the time
and axial terms this is \(H\), proving the first assertion.
The radial extensions and moment corrections preserve this axis
interval, so these equations concern the actual realized coefficients.

Now differentiate the full swirl sum, not a truncation:
\[
 B=\frac{2\sigma}{C}q^{-k}
 \left\{\varphi_0\left(\frac{\sigma}{\nu q},\eta\right)
       +\sum_{n\geq1}\chi(c_nq)q^{2nh}
          \varphi_n\left(\frac{\sigma}{\nu q},\eta\right)\right\}.
\]
Local finiteness justifies extracting its coefficient at every
preterminal point. The annular contributions vanish near the axis,
and the base streamfunction has no azimuthal curl component.
Using \eqref{nsr:finite-radial-closure} proves the second, independent
expression for the same full coefficient:
\begin{equation}\label{nsr:cutoff-b2}
 \boxed{\quad b_2(z,t)=\frac{q^{-2-h}H(\eta)
                   -\chi(c_1q)q^{-2+h}R(\eta)}{2\nu C}.\quad}
\end{equation}
Equating this with the full PDE identity \eqref{nsr:full-b2}
determines the force component exactly:
\begin{equation}\label{nsr:exact-force-curl}
 r_f(z,t)=\frac2C\{\chi(c_1q)-1\}q^{-2+h}R(\eta).
\end{equation}
The cutoff from Lemma~5.4, p.~57, is one on \([0,1/2]\).
Consequently \(r_f=0\) exactly whenever \(q\leq(2c_1)^{-1}\)
in this inner final-cutoff-one region. This is a proved cancellation
for the axial curl trace of the force, not a deletion of the actual
force and not an assertion of an unforced Navier--Stokes solution.
It is consistent with, and stronger for this coefficient than, the
flat estimate (10.9). The transition term in
\eqref{nsr:exact-force-curl} is retained outside that range.

\subsection{The fixed origin and a quantitative preterminal sign bound}

At \(z=0\), \(q=\tau\) and \(\eta=0\). Define the following
explicit constants without rescaling any field:
\begin{equation}\label{nsr:origin-constants}
 \begin{split}
 S_*&=j_0^2+\sigma_*^2,\qquad
 v=\varphi'(0)=-\frac{\Lambda j_0}{S_*},\\
 w&=\varphi''(0)
 =\frac{\Lambda(4+D)(j_0^2-\sigma_*^2)+\Lambda^2j_0^2}{S_*^2},\\
 K&=h-3+j_0v=h-3-\frac{\Lambda j_0^2}{S_*}<0,\qquad
 R_0=w-2(1+h).
 \end{split}
\end{equation}
To check the second derivative, \(H_*(0)=j_0\),
\(H_*'(0)=4+D\), \(L(0)=1\), and \(L'(0)=0\).
The quotient rule in \eqref{nsr:profiles} gives
\(\zeta_*'(0)=(4+D)(j_0^2-\sigma_*^2)/S_*^2\).
Since \(\varphi(0)=1\), its second derivative is
\(\Lambda\zeta_*'(0)+\Lambda^2\zeta_*(0)^2\), proving \(w\).
At zero \(Z_bv=v'\), but its second composition includes the
derivative of the coefficient \(2b\eta\); thus
\(R(0)=\varphi''(0)-2k\varphi(0)=R_0\).
Equation \eqref{nsr:full-b2} becomes the fully evaluated identity
\begin{equation}\label{nsr:origin-b2}
 \boxed{\quad
 b_2(0,1-\tau)=\frac{K\tau^{-2-h}-R_0\tau^{-2+h}}{2\nu C}
                            -\frac{r_f(0,1-\tau)}{4\nu}.
 \quad}
\end{equation}
In particular the factor \(\nu^{-1}\) is real coordinate information,
not an omitted choice of units.

Here is also a uniform version of the sign statement with a finite
threshold. Fix \(0<\eta_c<1\). For all sufficiently small \(q>0\),
the physical axis points
\[
 z=\sqrt\nu q^D\eta,\qquad t=1-q(1-\eta^2),\qquad |\eta|\leq\eta_c
\]
lie in the fixed localization neighborhood above. The profiles and
their parameter derivatives are smooth on this compact interval.
Let
\[
 m=(3-h-j_0)\min_{|\eta|\leq\eta_c}\varphi(\eta)>0,
 \qquad M=\max_{|\eta|\leq\eta_c}|R(\eta)|<\infty.
\]
Because \(L\leq1\), \eqref{nsr:H-sign} gives \(H\leq-m\).
Choose \(q_{\rm loc}>0\) to enforce both the source domain and
the final cutoffs equal to one. Explicitly, it suffices to take
\[
 q_{\rm loc}\leq\min\{q_{\rm source},\tau_0/2,(z_0/2)^{1/D}\},
\]
where \(q_{\rm source}\) is an inner source-domain bound, and
the original final cutoffs are one for \(\tau\leq\tau_0/2\)
and \(|y|\leq z_0/2\) near the radial axis. Indeed
\(\tau\leq q\) and \(|y|=q^D|\eta|\leq q^D\).
Put
\[
 q_*^{\rm sign}=\min\left\{1,q_{\rm loc},(2c_1)^{-1},
       \left(\frac{m}{2(M+1)}\right)^{1/(2h)}\right\}>0.
\]
For \(0<q<q_*^{\rm sign}\), multiply \eqref{nsr:full-b2} by
the positive number \(2\nu Cq^{2+h}\). The resulting exact
quantity is \(H-q^{2h}R-(C/2)q^{2+h}r_f\). The last term is
exactly zero by \eqref{nsr:exact-force-curl}; the middle term
has absolute value at most \(m/2\) by the displayed threshold.
Therefore
\begin{equation}\label{nsr:negative-bound}
 b_2(z,t)\leq-\frac{m}{4\nu C}q^{-2-h}<0
                 \quad(|\eta|\leq\eta_c,\ 0<q<q_*^{\rm sign}).
\end{equation}
No numerical value for the source-dependent threshold is claimed.
The formula retains its potentially small \(1/(2h)\) exponent.
The full force and every cutoff transition remain specified by
\eqref{nsr:extraction} and \eqref{nsr:exact-force-curl}.

\subsection{Exact arithmetic value, diffusion, and the remaining zero map}

Retain the full arithmetic transform, multiplier, and their inverse
from Section~\ref{sec:ns-hidden-mellin-diffusion}. With
\(\mathfrak B=\mathcal M\mathcal A_hB\), its exact regular value is
\begin{equation}\label{nsr:arithmetic-value}
 \boxed{\quad
 \mathfrak B(-2,z,t)
 =\frac{d_h}{2\nu C}
       \{q^{-2-h}H(\eta)-q^{-2+h}R(\eta)\}
                 -\frac{d_h}{4\nu}r_f(z,t),\qquad
 d_h=\frac{7(2^{h+2}-1)\zeta(3)}{32\pi^2}>0.
 \quad}
\end{equation}
The proof is \(\mathfrak B(-2)=W_h'(-2)b_2\) from
\eqref{ndm:all-coefficients}. Locally \(\widehat B(s)=
b_2/(s+2)+\text{holomorphic}\) and
\(W_h(s)=(s+2)(d_h+O(s+2))\); multiplication gives exactly
\eqref{nsr:arithmetic-value}. The classical functional equation
\cite[(25.4.2)]{apostolDLMF25} gives
\(\zeta'(-2)=-\zeta(3)/(4\pi^2)\), hence the displayed positive
constant. This is a regular continued Mellin value, not a convergent
unsubtracted integral at \(s=-2\). The explicit Taylor-subtracted
integral is \eqref{ndm:continuation}.

Equations \eqref{nsr:negative-bound} and \eqref{nsr:arithmetic-value}
prove a negative divergent regular Mellin value on the actual source
axis family. At the origin its leading coefficient is
\(d_hK/(2\nu C)\), with the exact subleading and force terms in
\eqref{nsr:origin-b2}. Yet the arithmetic trace's second derivative
is zero: \((\mathcal A_hB)''(0)=W_h(-2)B''(0)=0\).
The precise relation between these two statements is the zero-times-pole
product just calculated; the coefficient has not been destroyed.

For an explicit balance check at the origin, the other coefficients are
\[
 b_1=\frac2C\tau^{-1-h},\quad
 \partial_{zz}b_1=\frac{2R_0}{\nu C}\tau^{-2+h},\quad
 p_2=-\frac8C\tau^{-2-h},\quad
 \partial_zq_1=\frac2C(4+j_0v)\tau^{-2-h},\quad f_1=r_f.
\]
Substitution into \eqref{ndm:axis-balance} verifies every term.
Equivalently the full radial diffusion image has
\[
 (\mathscr L_h\mathcal A_hB)'(0)=4\kappa_hb_2.
\]
The arithmetic balance has right side
\[
 \frac{\kappa_h}{d_h}
      \{4\nu\mathfrak B(-2)-2\mathfrak P(-2)\}
 =\frac{2\kappa_h}{C}(5+h+j_0v)\tau^{-2-h}
     -\frac{2\kappa_hR_0}{C}\tau^{-2+h}-\kappa_h r_f.
\]
Its left side is exactly
\(\partial_tJ_B+\partial_zJ_Q-\nu\partial_{zz}J_B-J_F\).
Thus both nonlinear fluxes, axial diffusion, and force match the
growing residue; there is no scalar heat replacement or hidden compensation.

Finally this sign has a specific scope in the RH investigation.
It concerns the linear continued Mellin functional at the trivial
zero \(-2\), not the quadratic Weil form. The relation to zeros
elsewhere is the exact product
\(\mathfrak B(s)=\zeta(s)(1-2^{s-1})(1-2^{h-s})\widehat B(s)\).
In \(h<\Re s<1\), \(\widehat B\) is holomorphic by its original
integral, and both dilation factors are nonzero. At any point
\(s_0\) there, writing the nonzero analytic factors in their local
coordinates proves
\(\operatorname{ord}_{s_0}\mathfrak B
=\operatorname{ord}_{s_0}\zeta+\operatorname{ord}_{s_0}\widehat B\).
The actual \(B\) is nonzero by \eqref{nsr:axis-traces}, so Mellin
injectivity proves its transform is not identically zero; these
orders are well defined. The additional zero multiplicity is exactly
the actual fluid factor's order. Its nonvanishing at a proposed
nontrivial zeta zero, or a negative full Weil pairing, is not asserted
by the sign at \(-2\). This is the proved morphism and its full
zero accounting, not an inference of disconnection from different
presentations.
